% print(psi_forest_grid_tikz([4,5,4,4,3,2,1], upper_level=3, lower_level=3))
% Kupisch series [4,5,4,4,3,2,1]
  \centering
  {
  \newcommand{\nodea}{\node[draw=none, inner sep=2pt] (a) {}
  ;}
  \newcommand{\nodeb}{\node[draw=none, inner sep=2pt] (b) {}
  ;}
  \newcommand{\nodec}{\node[draw=none, inner sep=2pt] (c) {}
  ;}
  \newcommand{\noded}{\node[draw=none, inner sep=2pt] (d) {}
  ;}
  \newcommand{\nodee}{\node (e) {$\overline{10}$}
  ;}
  \newcommand{\nodef}{\node (f) {$\overline{9}$}
  ;}
  \newcommand{\nodeg}{\node (g) {$\overline{8}$}
  ;}
  \newcommand{\nodeh}{\node (h) {$\overline{7}$}
  ;}
  \newcommand{\nodei}{\node (i) {$\overline{6}$}
  ;}
  \newcommand{\nodej}{\node (j) {$\overline{5}$}
  ;}
  \newcommand{\nodek}{\node (k) {$\overline{4}$}
  ;}
  \newcommand{\nodel}{\node (l) {$\overline{3}$}
  ;}
  \newcommand{\nodem}{\node (m) {$\overline{2}$}
  ;}
  \newcommand{\noden}{\node (n) {$\overline{1}$}
  ;}
  \newcommand{\nodeo}{\node (o) {$0$}
  ;}
  \newcommand{\nodep}{\node (p) {$1$}
  ;}
  \newcommand{\nodeq}{\node (q) {$2$}
  ;}
  \newcommand{\noder}{\node (r) {$3$}
  ;}
  \newcommand{\nodes}{\node (s) {$4$}
  ;}
  \newcommand{\nodet}{\node (t) {$5$}
  ;}
  \newcommand{\nodeu}{\node (u) {$6$}
  ;}
  \newcommand{\nodev}{\node (v) {$7$}
  ;}
  \newcommand{\nodew}{\node (w) {$8$}
  ;}
  \newcommand{\nodex}{\node (x) {$9$}
  ;}
  \newcommand{\nodey}{\node (y) {$10$}
  ;}
  \newcommand{\nodez}{\node (z) {$11$}
  ;}
  \newcommand{\nodeaa}{\node (aa) {$12$}
  ;}
  \newcommand{\nodeab}{\node (ab) {$13$}
  ;}
  \newcommand{\nodeac}{\node[draw=none, inner sep=2pt] (ac) {}
  ;}
  \newcommand{\nodead}{\node[draw=none, inner sep=2pt] (ad) {}
  ;}
  \newcommand{\nodeae}{\node[draw=none, inner sep=2pt] (ae) {}
  ;}
  \newcommand{\nodeaf}{\node[draw=none, inner sep=2pt] (af) {}
  ;}
  \begin{tikzpicture}[auto, font=\footnotesize]
  \matrix[column sep=.3cm, row sep=.5cm,ampersand replacement=\&]{
    \nodea  \&  \nodeb  \&  \nodec  \&  \noded  \\
    \nodee  \&  \nodef  \&  \nodeg  \&           \\
    \nodeh  \&  \nodei  \&  \nodej  \&  \nodek  \\
    \nodel  \&  \nodem  \&  \noden  \&           \\
    \nodeo  \&  \nodep  \&  \nodeq  \&  \noder  \\
    \nodes  \&  \nodet  \&  \nodeu  \&           \\
    \nodev  \&  \nodew  \&  \nodex  \&  \nodey  \\
    \nodez  \&  \nodeaa  \&  \nodeab  \&           \\
    \nodeac  \&  \nodead  \&  \nodeae  \&  \nodeaf  \\
  };

  \path[ultra thick]   (g) edge (h)
  	  (g) edge (k)
  	  (h) edge (l)
  	  (h) edge (m)
  	  (j) edge (n)
  	  (n) edge (o)
  	  (n) edge (p)
  	  (n) edge (q)
  	  (n) edge (r)
  	  (o) edge (s)
  	  (o) edge (t)
  	  (u) edge (v)
  	  (u) edge (w)
  	  (u) edge (x)
  	  (u) edge (y)
  	  (x) edge (ab);
  \path[ultra thick, dotted]   (a) edge (e)
  	  (a) edge (f)
  	  (c) edge (g)
  	  (ab) edge (ac)
  	  (ab) edge (ad)
  	  (ab) edge (ae)
  	  (ab) edge (af);
  % Colour the edges of the three paths in the colours of their endpoint vertices
  \path[ultra thick, draw=green]  (v) edge (z)  (v) edge (aa);
  \path[ultra thick, draw=red]    (g) edge (i)  (g) edge (j);
  \path[ultra thick, draw=blue]   (q) edge (u);
  \begin{pgfonlayer}{background}
  % Highlight the triplets of vertices in the colours of their modules
  \foreach \v in {i,j,m}{\node[circle, fill=red, opacity=0.35, inner sep=4pt] at (\v) {};}
  \foreach \v in {q,r,u}{\node[circle, fill=blue, opacity=0.35, inner sep=4pt] at (\v) {};}
  \foreach \v in {z,aa,w}{\node[circle, fill=green, opacity=0.35, inner sep=4pt] at (\v) {};}
  \end{pgfonlayer}
  \end{tikzpicture}}
